% ============================================================================
% Chapter 2 — Materials and Methods (shared general methods)
% ============================================================================
\chapter{Materials and Methods}
\label{ch:methods}

\section{Reagents and solvents}
\label{sec:reagents}
All reagents were purchased from commercial suppliers and used as received
unless otherwise stated. Dry solvents (THF, DCM, toluene) were obtained from a
solvent purification system. Column chromatography was performed on silica gel
(100--200 mesh).

\section{Instrumentation}
\label{sec:instruments}
NMR: Bruker Avance \SI{400}{\mega\hertz} and \SI{500}{\mega\hertz}
spectrometers; chemical shifts in ppm relative to TMS; coupling constants in Hz.
HRMS: Bruker microTOF-Q II (ESI). IR: PerkinElmer Spectrum Two (ATR). UV--vis:
Shimadzu UV-2600. Fluorescence: Horiba Fluorolog-3. SCXRD: Bruker D8 Venture.
CV: CH Instruments 600E. Computational: Gaussian 16 / ORCA 5.0.

\section{General experimental procedures}
\label{sec:genproc}
\subsection{Standard characterisation protocol}
\label{sec:charproto}
Every new compound was characterised by \ce{^1H} and \ce{^{13}C}\{\ce{^1H}\} NMR,
HRMS, IR and melting point; purity was confirmed by elemental analysis where
required. Crystals suitable for SCXRD were deposited at the CCDC (numbers
quoted per chapter); checkCIF reports are compiled in Appendix A.

\section{Computational methods}
\label{sec:compmethods}
Geometry optimisations were performed at the \texttt{PBE0-D3BJ/def2-TZVP} level
using Gaussian 16 (Revision C.01). Solvation effects were included via the
IEFPCM model (THF). All stationary points were verified by frequency analysis.
Coordinates are listed in Appendix C.
